% ============================================================
%  Chapitre 4 — Modélisation : architecture du système RAG
% ============================================================
\chapter{Modélisation : architecture du système RAG}
\label{ch:modelisation}

\section{Introduction}

Ce chapitre présente le cœur technique du projet. Comme souligné au
chapitre~\ref{ch:contexte}, la phase de modélisation ne consiste pas ici à
entraîner un modèle, mais à \textbf{assembler et paramétrer une chaîne de
traitement} dont chaque maillon influence la qualité finale. Nous détaillons
successivement les quatre composants de cette chaîne : la vectorisation des
segments, leur indexation dans une base vectorielle, la recherche sémantique, et
la génération contrainte de la réponse. Nous exposons enfin la manière dont
LangChain orchestre l'ensemble.

\section{Architecture générale}

Le système repose sur une séparation stricte entre deux phases, illustrée à la
figure~\ref{fig:archi}.

La \textbf{phase hors ligne} (ingestion) est exécutée une fois, puis à chaque
évolution du corpus. Elle transforme les documents bruts en un index vectoriel
persistant. Son coût est élevé mais amorti.

La \textbf{phase en ligne} (interrogation) est exécutée à chaque question. Elle
ne touche jamais aux fichiers d'origine : elle interroge l'index, construit une
invite et appelle le modèle de langue.

Ce découplage est une décision d'architecture importante. Il permet de
ré-indexer le corpus sans interrompre le service, et de faire évoluer la
stratégie de recherche ou l'invite système sans avoir à re-vectoriser les
2\,929 segments --- opération qui prend plusieurs minutes sur CPU.

\begin{figure}[htbp]
\centering
\begin{tikzpicture}[
  font=\small,
  box/.style={draw=ink,line width=0.9pt,rounded corners=4pt,fill=white,
              text width=3.05cm,align=center,minimum height=0.95cm},
  acc/.style={draw=accent,line width=1.1pt,rounded corners=4pt,fill=soft,
              text width=3.05cm,align=center,minimum height=0.95cm},
  store/.style={draw=accent,line width=1.2pt,cylinder,shape border rotate=90,
              aspect=0.22,fill=soft,text width=2.5cm,align=center,
              minimum height=1.2cm,inner sep=4pt},
  fl/.style={-{Stealth[length=2.4mm]},line width=0.9pt,color=muted},
  lbl/.style={font=\scriptsize\itshape,color=muted}
]

% ---------- phase offline (colonne de gauche) ----------
\node[box] (docs) {Documents bruts\\{\scriptsize PDF, DOCX, TXT, MD, CSV}};
\node[box,below=0.75cm of docs] (load) {\textbf{Chargement}\\{\scriptsize extraction texte\\+ métadonnées}};
\node[box,below=0.75cm of load] (chunk) {\textbf{Découpage}\\{\scriptsize 1000 car. / 150 chev.}};
\node[acc,below=0.75cm of chunk] (embed) {\textbf{Vectorisation}\\{\scriptsize all-MiniLM-L6-v2\\384 dimensions}};
\node[store,below=0.85cm of embed] (chroma) {\textbf{ChromaDB}\\{\scriptsize 2\,929 segments}};

\draw[fl] (docs) -- (load);
\draw[fl] (load) -- (chunk);
\draw[fl] (chunk) -- (embed);
\draw[fl] (embed) -- (chroma);

% ---------- phase online (colonne de droite) ----------
\node[box,right=3.4cm of docs] (quest) {\textbf{Question}\\{\scriptsize langage naturel}};
\node[acc,below=0.75cm of quest] (qemb) {\textbf{Vectorisation}\\{\scriptsize même modèle}};
\node[acc,below=0.75cm of qemb] (search) {\textbf{Recherche}\\{\scriptsize similarité cosinus\\top-$k$ = 4}};
\node[box,below=0.75cm of search] (prompt) {\textbf{Invite}\\{\scriptsize contexte + question\\+ consignes}};
\node[acc,below=0.75cm of prompt] (llm) {\textbf{Génération}\\{\scriptsize Ollama / Mistral 7B\\$T$ = 0,1}};
\node[box,below=0.75cm of llm] (rep) {\textbf{Réponse sourcée}\\{\scriptsize + extraits + scores}};

\draw[fl] (quest) -- (qemb);
\draw[fl] (qemb) -- (search);
\draw[fl] (search) -- (prompt);
\draw[fl] (prompt) -- (llm);
\draw[fl] (llm) -- (rep);

% ---------- lien entre les deux phases ----------
\draw[fl,color=accent,line width=1.1pt] (chroma.east) -- ++(0.55,0)
      |- node[lbl,pos=0.75,above]{top-$k$ extraits} (search.west);

% ---------- étiquettes de phase ----------
\node[lbl,above=0.18cm of docs]  {\textbf{Phase hors ligne} --- ingestion};
\node[lbl,above=0.18cm of quest] {\textbf{Phase en ligne} --- interrogation};

\end{tikzpicture}
\caption{Architecture du système RAG. La phase hors ligne construit l'index
vectoriel ; la phase en ligne l'interroge à chaque question. Le modèle de
vectorisation est nécessairement le même dans les deux phases.}
\label{fig:archi}
\end{figure}

Un point mérite d'être souligné sur cette figure : le modèle de vectorisation
apparaît \emph{deux fois}, et il doit impérativement s'agir du même. Les vecteurs
des segments et celui de la question doivent appartenir au même espace
sémantique pour que leur comparaison ait un sens. Changer de modèle
d'\eng{embeddings} impose donc de reconstruire intégralement l'index --- une
contrainte opérationnelle dont il faut tenir compte avant tout changement.

\section{Vectorisation sémantique}

\subsection{Principe}

La vectorisation, ou \eng{embedding}, consiste à projeter un texte dans un espace
vectoriel de dimension fixe, de telle sorte que la proximité géométrique entre
deux vecteurs traduise la proximité de sens entre les deux textes. C'est ce
mécanisme qui résout le fossé lexical identifié au chapitre~\ref{ch:metier} :
« congés » et « absences autorisées » n'ont aucun caractère en commun, mais
leurs vecteurs sont voisins.

Formellement, le modèle réalise une application
\[
  E : \mathcal{T} \longrightarrow \mathbb{R}^{d}
\]
où $\mathcal{T}$ est l'ensemble des textes et $d$ la dimension de l'espace
d'arrivée. Pour le modèle retenu, $d = 384$.

\subsection{Choix du modèle}

Le modèle \texttt{sentence-transformers/all-MiniLM-L6-v2} a été retenu. Il
s'agit d'une architecture Transformer à 6 couches, distillée depuis un modèle
plus volumineux, comportant environ 22,7 millions de paramètres et produisant des
vecteurs de dimension 384.

Ce choix résulte d'un compromis explicite entre trois exigences.
La \textbf{gratuité et la confidentialité} d'abord : le modèle s'exécute
localement, sans appel réseau, contrairement à l'API
\texttt{text-embedding-ada-002} d'OpenAI mentionnée dans le sujet.
La \textbf{légèreté} ensuite : avec 384 dimensions contre 1\,536 pour le modèle
d'OpenAI, l'index occupe quatre fois moins d'espace et la recherche est plus
rapide ; le modèle tourne confortablement sur CPU.
La \textbf{qualité} enfin : ce modèle constitue une référence établie pour la
recherche sémantique en anglais, langue dans laquelle le corpus est rédigé.

\begin{encadre}
\textbf{Limite identifiée.} Le modèle \texttt{all-MiniLM-L6-v2} est
majoritairement entraîné sur de l'anglais. Or les questions des utilisateurs sont
formulées en français. Cette asymétrie dégrade la qualité de l'appariement
lorsque la question et le document n'utilisent pas la même langue. Le passage à
\texttt{paraphrase-multilingual-MiniLM-L12-v2}, entraîné sur plus de cinquante
langues, constitue la première piste d'amélioration identifiée
(\emph{cf.} chapitre~\ref{ch:evaluation}). Ce changement impose toutefois une
ré-indexation complète du corpus.
\end{encadre}

\subsection{Normalisation}

Les vecteurs sont normalisés à la production ($\lVert v \rVert_2 = 1$). Cette
normalisation présente un avantage pratique : la similarité cosinus se réduit
alors au simple produit scalaire,
\[
  \cos(u, v) = \frac{u \cdot v}{\lVert u \rVert\,\lVert v \rVert} = u \cdot v,
\]
ce qui accélère la recherche et rend les scores directement comparables entre
requêtes.

\begin{lstlisting}[style=python,caption={Instanciation du modèle de vectorisation (\texttt{src/ingestion/indexer.py})}]
def get_embeddings() -> HuggingFaceEmbeddings:
    """Instancie paresseusement le modele d'embeddings (charge une fois, reutilise)."""
    global _embeddings_instance
    if _embeddings_instance is None:
        _embeddings_instance = HuggingFaceEmbeddings(
            model_name=settings.embedding_model,
            encode_kwargs={"normalize_embeddings": True},
        )
    return _embeddings_instance
\end{lstlisting}

Le chargement est \emph{paresseux} et mis en cache dans une variable de module :
le modèle, qui pèse environ 90~Mo, n'est chargé qu'au premier appel puis
réutilisé. Sans cette précaution, chaque requête HTTP rechargerait le modèle
depuis le disque, ajoutant plusieurs secondes de latence.

\section{Indexation vectorielle}

\subsection{ChromaDB}

Les vecteurs produits sont persistés dans une collection ChromaDB nommée
\texttt{company\_documents}, stockée localement dans
\texttt{data/processed/chroma\_db/}. ChromaDB a été préférée à Pinecone et Qdrant
pour trois raisons : elle fonctionne intégralement en local sans serveur ni clé
d'API, elle s'intègre nativement à LangChain, et ses performances sont largement
suffisantes à l'échelle de plusieurs centaines de milliers de segments.

En interne, ChromaDB indexe les vecteurs à l'aide d'un graphe \textbf{HNSW}
(\eng{Hierarchical Navigable Small World}), une structure de recherche
approximative de plus proches voisins. Elle évite le parcours exhaustif de la
collection : la complexité de recherche passe de $O(n)$ à
$O(\log n)$ environ, au prix d'une approximation en pratique négligeable.

\subsection{Indexation par lots}

L'insertion s'effectue par lots de 128 segments plutôt que segment par segment.
Ce choix réduit sensiblement la durée totale d'indexation, la vectorisation
bénéficiant du traitement vectoriel par lot, et il limite l'empreinte mémoire
en évitant de maintenir l'intégralité des vecteurs en RAM.

\begin{lstlisting}[style=python,caption={Indexation par lots (\texttt{src/ingestion/indexer.py})}]
def index_documents(chunks: list[Document], batch_size: int = 128) -> int:
    """Vectorise et persiste les segments dans ChromaDB."""
    if not chunks:
        return 0
    vectorstore = get_vectorstore()
    total = 0
    for i in range(0, len(chunks), batch_size):
        batch = chunks[i : i + batch_size]
        vectorstore.add_documents(batch)
        total += len(batch)
    return total
\end{lstlisting}

L'état final de l'index est présenté au tableau~\ref{tab:index}.

\begin{table}[htbp]
\centering
\caption{Caractéristiques de l'index vectoriel construit}
\label{tab:index}
\small
\begin{tabularx}{\textwidth}{@{}p{5.4cm}X@{}}
\toprule
\textbf{Paramètre} & \textbf{Valeur} \\
\midrule
Base vectorielle & ChromaDB (persistée sur disque local) \\
Nom de la collection & \texttt{company\_documents} \\
Nombre de segments indexés & 2\,929 \\
Modèle de vectorisation & \texttt{all-MiniLM-L6-v2} \\
Dimension des vecteurs & 384 \\
Métrique de similarité & Cosinus (vecteurs normalisés) \\
Structure d'index & HNSW \\
Taille sur disque & environ 18~Mo \\
Taille de lot à l'insertion & 128 segments \\
\bottomrule
\end{tabularx}
\end{table}

\section{Recherche sémantique}

À l'interrogation, la question est vectorisée par le même modèle, puis comparée à
l'ensemble des vecteurs de l'index. Les $k$ segments les plus proches sont
retournés, accompagnés de leur score de pertinence.

Le paramètre $k$ a été fixé à \textbf{4}. Sa valeur résulte d'un arbitrage :
un $k$ trop faible risque d'omettre l'information nécessaire (dégradant le
\eng{context recall}) ; un $k$ trop élevé introduit des segments hors sujet qui
diluent le contexte, dégradent le \eng{context precision} et augmentent le
risque que le modèle s'appuie sur un passage non pertinent. La valeur $k = 4$
constitue par ailleurs un compromis de latence : chaque segment
supplémentaire allonge l'invite et donc le temps de génération.

\begin{lstlisting}[style=python,caption={Recherche avec scores (\texttt{src/rag/retriever.py})}]
def retrieve_with_scores(query: str, k: int | None = None):
    """Recupere les segments avec leur score de pertinence."""
    vectorstore = get_vectorstore()
    return vectorstore.similarity_search_with_relevance_scores(
        query, k=k or settings.top_k
    )
\end{lstlisting}

L'exposition des scores n'est pas un simple élément de confort de mise au point.
Elle remplit deux fonctions dans le système final : elle alimente
l'affichage de traçabilité présenté à l'utilisateur, et elle sert de signal pour
le garde-fou décrit à la section~\ref{sec:garde-fous}.

\section{Génération contrainte}

\subsection{Le modèle de langue}

La génération est assurée par \textbf{Mistral 7B}, servi localement par
\textbf{Ollama}. Le modèle, quantifié, occupe environ 4,4~Go sur disque et
s'exécute sur CPU. Ollama expose une API HTTP locale sur le port 11434, à
laquelle LangChain se connecte via le connecteur \texttt{ChatOllama}.

La \textbf{température} est fixée à \textbf{0,1}, valeur volontairement très
basse. La température contrôle l'aléa de l'échantillonnage des \eng{tokens} :
plus elle est élevée, plus le modèle produit un texte varié et créatif. Or la
créativité est exactement ce qu'il faut proscrire ici. Nous attendons du modèle
qu'il restitue fidèlement le contenu du contexte fourni, non qu'il le reformule
librement. Une température quasi nulle rend en outre les réponses largement
reproductibles, ce qui est indispensable pour que l'évaluation du
chapitre~\ref{ch:evaluation} soit stable d'une exécution à l'autre.

\subsection{L'invite système : premier levier anti-hallucination}

L'invite système constitue le principal levier de réduction des hallucinations
au niveau de la génération. Elle a été formulée pour contraindre explicitement
le modèle sur quatre points : la restriction au contexte fourni, l'obligation
d'avouer son ignorance, l'interdiction d'inventer des faits, et la citation des
sources.

\begin{lstlisting}[style=plain,caption={Invite système (\texttt{src/rag/pipeline.py})}]
Tu es un assistant interne d'entreprise. Reponds UNIQUEMENT a partir du
contexte fourni ci-dessous, extrait des documents internes de l'entreprise.

Regles strictes :
- Si la reponse ne se trouve pas dans le contexte, reponds explicitement :
  "Je ne trouve pas cette information dans les documents disponibles."
- Ne jamais inventer de faits, chiffres ou procedures qui ne sont pas
  dans le contexte.
- Cite les documents sources utilises (nom de fichier) a la fin de ta reponse.
- Reponds de facon concise et professionnelle, en francais, sauf si la
  question est posee en anglais.

Contexte :
{context}

Question : {question}

Reponse :
\end{lstlisting}

Deux détails de formulation méritent d'être relevés. D'une part, la consigne
d'ignorance fournit au modèle une \emph{formulation exacte} à reprendre plutôt
qu'une simple interdiction : lui donner une sortie de secours explicite est
nettement plus efficace que de lui interdire d'inventer sans lui indiquer quoi
faire à la place. D'autre part, chaque extrait injecté dans le contexte est
préfixé du nom de son fichier d'origine, ce qui permet au modèle de citer
correctement ses sources :

\begin{lstlisting}[style=python,caption={Formatage du contexte avec attribution de source}]
def _format_docs(docs) -> str:
    return "\n\n---\n\n".join(
        f"[Source: {d.metadata.get('filename', 'unknown')}]\n{d.page_content}"
        for d in docs
    )
\end{lstlisting}

\subsection{Garde-fou sur le contexte vide}
\label{sec:garde-fous}

L'invite ne suffit pas. Un second garde-fou, de nature programmatique,
court-circuite entièrement le modèle lorsque la recherche ne remonte aucun
segment. Dans ce cas, la réponse d'ignorance est retournée directement, sans
appel au LLM.

\begin{lstlisting}[style=python,caption={Garde-fou programmatique (\texttt{src/rag/pipeline.py})}]
scored_docs = retrieve_with_scores(question, k=k)

if not scored_docs:
    return RAGResponse(
        answer="Je ne trouve pas cette information dans les documents disponibles.",
        sources=[],
        context_found=False,
    )
\end{lstlisting}

Ce garde-fou est plus robuste que l'invite, car il ne dépend pas de la
docilité du modèle : un LLM peut ignorer une consigne, il ne peut pas contourner
une instruction conditionnelle. C'est un principe de conception qu'il convient
de généraliser : \emph{tout ce qui peut être garanti par du code ne doit pas être
délégué à l'invite}.

\section{Orchestration avec LangChain}

L'ensemble est assemblé au moyen du langage d'expression de chaînes de LangChain
(\eng{LangChain Expression Language}), qui permet de composer les étapes à
l'aide de l'opérateur \texttt{|}.

\begin{lstlisting}[style=python,caption={Composition de la chaîne RAG (\texttt{src/rag/pipeline.py})}]
chain = (
    {"context": lambda x: _format_docs(docs), "question": RunnablePassthrough()}
    | _prompt
    | get_llm()
    | StrOutputParser()
)

answer = chain.invoke(question)
\end{lstlisting}

La chaîne se lit de gauche à droite : construction du dictionnaire de variables,
application du gabarit d'invite, appel du modèle, extraction de la chaîne de
caractères de sortie. Cette écriture déclarative rend la structure du pipeline
immédiatement lisible et facilite le remplacement d'un maillon --- changer de
modèle de langue ne demande de modifier qu'un seul appel.

La réponse finale agrège le texte généré et les métadonnées de traçabilité :

\begin{lstlisting}[style=python,caption={Construction de la réponse sourcée}]
sources = [
    {
        "filename": d.metadata.get("filename"),
        "chunk_id": d.metadata.get("chunk_id"),
        "relevance_score": round(float(score), 3),
    }
    for d, score in scored_docs
]

return RAGResponse(answer=answer, sources=sources, context_found=True)
\end{lstlisting}

\section{Synthèse des paramètres}

Le tableau~\ref{tab:hyperparams} récapitule l'ensemble des paramètres du
pipeline. Tous sont exposés par variables d'environnement, ce qui permet de les
faire varier sans modification du code --- condition nécessaire à l'itération
mesurée décrite au chapitre suivant.

\begin{table}[htbp]
\centering
\caption{Paramètres du pipeline RAG}
\label{tab:hyperparams}
\small
\begin{tabularx}{\textwidth}{@{}p{3.6cm}p{3.4cm}X@{}}
\toprule
\textbf{Paramètre} & \textbf{Valeur} & \textbf{Justification} \\
\midrule
\texttt{CHUNK\_SIZE} & 1000 caractères &
Assez large pour conserver un document transactionnel entier, assez court pour
rester sous la limite du modèle d'\eng{embeddings} \\
\texttt{CHUNK\_OVERLAP} & 150 caractères &
Évite qu'une information soit coupée à la frontière de deux segments \\
\texttt{TOP\_K} & 4 &
Compromis entre \eng{recall} (ne rien manquer) et \eng{precision} (ne pas diluer
le contexte) \\
\texttt{EMBEDDING\_MODEL} & \texttt{all-MiniLM-L6-v2} &
Local, léger (384 dim.), gratuit \\
\texttt{OLLAMA\_MODEL} & \texttt{mistral} &
7 milliards de paramètres, exécutable sur CPU, licence ouverte \\
\texttt{OLLAMA\_TEMPERATURE} & 0,1 &
Minimise la créativité et rend les réponses reproductibles \\
\bottomrule
\end{tabularx}
\end{table}

\section{Conclusion}

Ce chapitre a détaillé l'architecture du système : vectorisation par
\texttt{all-MiniLM-L6-v2}, indexation HNSW dans ChromaDB, recherche par
similarité cosinus à $k = 4$, et génération contrainte par Mistral 7B à
température 0,1. Nous avons insisté sur les deux leviers anti-hallucination mis
en place --- une invite système explicite et un garde-fou programmatique sur le
contexte vide --- ainsi que sur le principe qui les sous-tend : privilégier la
garantie par le code sur la consigne en langage naturel. Reste à vérifier que ces
leviers sont effectivement efficaces. C'est l'objet du chapitre suivant.
